\documentclass[10pt,a4paper]{amsart}
\usepackage[margin=19mm]{geometry}
\usepackage{amsmath,amssymb,mathtools}
\usepackage[scr=rsfs,cal=euler]{mathalfa}
\usepackage{xfrac}
\usepackage[unicode,hidelinks,bookmarks=false]{hyperref}
\newcommand{\ie}{i.e.\ }
\newcommand{\cf}{cf.\ }
\newcommand{\resp}{respectively\ }
\newcommand{\Subsection}{\subsection}
\providecommand{\nobreakdash}{\nobreak\hbox{-}}
\setlength{\emergencystretch}{2em}
\multlinegap0pt
\usepackage{bm}
\usepackage{bbm}
\usepackage{dsfont}
\usepackage{xspace}
\usepackage{cancel}
\usepackage{mathtools}
\usepackage{amsthm}
\makeatletter
\@ifundefined{theorem}{\newtheorem{theorem}{Theorem}}{}
\makeatother
\title[Global iterative methods for sparse approximate inverses of ]{Global iterative methods for sparse approximate inverses of symmetric positive definite matrices}
\author{Nicolas Venkovic, Hartwig Anzt}
\date{Source: arXiv:2511.09753v2 (CC BY 4.0)}
\begin{document}
\maketitle

\noindent\emph{Source and licence.} Global iterative methods for sparse approximate inverses of symmetric positive definite matrices. Version arXiv:2511.09753v2 (CC BY 4.0), distributed under the Creative Commons Attribution 4.0 International licence, as recorded in the arXiv licence metadata for this exact version and shown on the abstract page https://arxiv.org/abs/2511.09753v2. Full rights notice: licenses/arxiv-2511.09753-CC-BY-4.0.txt.\par\smallskip

\noindent\textit{Editorial scope.} Counted unit: the Theorem whose source label is \texttt{theo:mr-convergence} in the article (source0.tex lines 250--253, source label \texttt{theo:mr-convergence}), delivered together with the article's own complete original proof of it (source0.tex lines 255--322, one \texttt{proof} environment of 68 lines). No mathematics is written, repaired or re-derived: statement and proof are the article's own text line for line. Required context retained uncounted: none. Every definition, display and label of those lines is retained unchanged. Omitted from this package and not redistributed: 1--249, 254--254, 323--1635. Nothing inside a retained excerpt is omitted or altered: every retained body is delivered exactly as the article's lines are written, so no omission receipt of any kind accompanies this package. Citations: the retained excerpts carry no citation command at all, so no bibliography entry of the article is transcribed and no reference is redistributed. Reference adaptation: none was needed and none was made; every reference inside the delivered unit resolves inside it, so no unresolved reference or citation is produced. Printed slips kept literal: no printed word, symbol, hypothesis or number of the counted statement or of its proof was altered. Layout and class: the article's own class is \texttt{ima-authoring-template} with packages including cancel, relsize, stmaryrd, comment, caption, tikz, booktabs; the wrapper is the lane's pinned \texttt{amsart} editorial wrapper at 10pt on A4, which restates the theorem-like environments the retained text needs and adds the article's own macro definitions that the retained text needs (none), with extra wrapper packages (bm, bbm, dsfont, xspace, cancel, mathtools). The delivered numbering is the unit's own. Assets: no retained span contains a figure environment, a graphics-inclusion command, a PDF-page inclusion, a table or a TikZ picture; no image file is copied into this package, the package declares no asset dependency and no image file is redistributed with it. Licence: version arXiv:2511.09753v2 of this article is distributed under the Creative Commons Attribution 4.0 International licence, as recorded in the arXiv licence metadata for this exact version and shown on the abstract page https://arxiv.org/abs/2511.09753v2. A full rights notice is delivered at licenses/arxiv-2511.09753-CC-BY-4.0.txt. Characters: every retained excerpt is pure ASCII, so the delivered engine, XeLaTeX with its default Latin Modern text font, reports no missing character.
\begin{theorem}[Convergence of MR iterates]\label{theo:mr-convergence}
Irrespective of the initial guess $M_0$, the MR iterates $M_0,M_1,\dots$ converge to the unique solution $M=A^{-1}$ such that $AM=I_n$.
They do so linearly at a rate $(1-\lambda_{\min}(A)^2/\operatorname{tr}(A^2))^{1/2}$.
\end{theorem}

\begin{proof}[Proof of Theo.~\ref{theo:mr-convergence}]
For every new MR iterate $M_{i+1}$, the residual norm is given by
\begin{align*}
\|R_{i+1}\|_F^2
=(R_{i+1},R_{i+1})_F
=(R_{i+1},R_i)_F-\alpha_i(R_{i+1},AR_i)_F\\
\end{align*}
where, by definition, we have the orthogonality condition $R_{i+1}\perp A\operatorname{span}\{R_i\}$ so that
\begin{align*}
\|R_{i+1}\|_F^2
=(R_{i+1},R_i)_F
\end{align*}
and
\begin{align}
\|R_{i+1}\|_F^2
=(R_{i+1},R_i)_F
=(R_i-\alpha_i AR_i,R_i)_F
=\|R_i\|_F^2-\alpha_i(AR_i,R_i)_F.
\end{align}
Using the fact that 
\begin{align}
\alpha_i=\frac{(AR_i,R_i)_F}{(AR_i,AR_i)_F},
\end{align}
we get
\begin{align}
\|R_{i+1}\|_F^2
=&\,\|R_i\|_F^2-\frac{(AR_i,R_i)_F}{(AR_i,AR_i)_F}(AR_i,R_i)_F\\
=&\,\|R_i\|_F^2\left(1-\frac{(AR_i,R_i)_F}{(AR_i,AR_i)_F}\frac{(AR_i,R_i)_F}{\|R_i\|_F^2}\right)\\
=&\,\|R_i\|_F^2\left(1-\frac{(AR_i,R_i)_F^2}{\|AR_i\|_F^2}\frac{\|R_i\|_F^2}{\|R_i\|_F^4}\right)
\end{align}
where, by consistency of the Frobenius norm, we have $\|AR_i\|_F\leq\|A\|_F\|R_i\|_F$ such that 
\begin{align}
\|R_{i+1}\|_F^2
\leq&\,\|R_i\|_F^2\left(1-\frac{(AR_i,R_i)_F^2}{\|A\|_F^2\|R_i\|_F^2}\frac{\|R_i\|_F^2}{\|R_i\|_F^4}\right)
\end{align}
eventually leading to
\begin{align}\label{eq:1000}
\|R_{i+1}\|_F^2
\leq&\,\|R_i\|_F^2\left(1-\frac{(AR_i,R_i)_F^2}{\|A\|_F^2(R_i,R_i)_F^2}\right).
\end{align}
Let us then reformulate the quotient as
\begin{align}
\frac{(AR_i,R_i)_F}{(R_i,R_i)_F}
=\frac{\operatorname{vec}(AR_i)^T\operatorname{vec}(R_i)}{\operatorname{vec}(R_i)^T\operatorname{vec}(R_i)}
\end{align}
where $\operatorname{vec}(AR_i)=(I_n\otimes A)\operatorname{vec}(R_i)$ so that
\begin{align}\label{eq:1010}
\frac{(AR_i,R_i)_F}{(R_i,R_i)_F}
=\frac{((I_n\otimes A)\operatorname{vec}(R_i))^T\operatorname{vec}(R_i)}{\operatorname{vec}(R_i)^T\operatorname{vec}(R_i)}
=\frac{\operatorname{vec}(R_i)^T(I_n\otimes A)\operatorname{vec}(R_i)}{\operatorname{vec}(R_i)^T\operatorname{vec}(R_i)}.
\end{align}
The spectrum of $I_n\otimes A$ is the same as that of $A$.
Thus, since $A$ is SPD, we have 
\begin{align}\label{eq:1020}
\frac{\operatorname{vec}(R_i)^T(I_n\otimes A)\operatorname{vec}(R_i)}{\operatorname{vec}(R_i)^T\operatorname{vec}(R_i)}\geq \lambda_{\min}(I_n\otimes A)=\lambda_{\min}(A)>0.
\end{align}
Combining Eqs.~\eqref{eq:1000}, \eqref{eq:1010} and \eqref{eq:1020} with the fact that $\|A\|_F^2=\operatorname{tr}(A^2)$, we obtain
\begin{align}
\|R_{i+1}\|_F^2\leq\|R_i\|^2\left(1-\frac{\lambda_{\min}(A)^2}{\operatorname{tr}(A^2)}\right)
\end{align}
and thus
\begin{align}
\|R_{i+1}\|_F\leq C\|R_i\|
\;\text{ with }\;
C:=\left(1-\frac{\lambda_{\min}(A)^2}{\operatorname{tr}(A^2)}\right)^{1/2}<1. 
\end{align}

\end{proof}

\end{document}
